\documentclass[aspectratio=169,10pt]{beamer}
\usetheme[numbering=fraction,progressbar=frametitle]{metropolis}
\usepackage{fontspec}
\setsansfont{Latin Modern Sans}
\usepackage{amsmath,amssymb,mathtools,graphicx,booktabs}
\usepackage{polyglossia}
\setdefaultlanguage{english}
\setbeamertemplate{navigation symbols}{}
\setbeamerfont{frametitle}{size=\large}
\setbeamerfont{normal text}{size=\normalsize}
\setlength{\parskip}{.45em}
\hypersetup{pdfauthor={},pdfsubject={Lecture slides},pdfkeywords={},pdftitle={Prime numbers and Fermat's little theorem}}
\providecommand{\N}{\mathbb{N}}
\providecommand{\Z}{\mathbb{Z}}
\providecommand{\D}{\mathbb{D}}
\providecommand{\Q}{\mathbb{Q}}
\providecommand{\R}{\mathbb{R}}
\providecommand{\C}{\mathbb{C}}
\providecommand{\vect}[1]{\overrightarrow{#1}}
\providecommand{\vi}{\vec{\imath}}
\providecommand{\vj}{\vec{\jmath}}
\providecommand{\vk}{\vec{k}}
\providecommand{\norm}[1]{\left\lVert #1\right\rVert}
\providecommand{\vc}[2]{\begin{pmatrix}#1\\#2\end{pmatrix}}
\providecommand{\deter}[2]{\det\left(#1,#2\right)}
\providecommand{\abs}[1]{\left\lvert #1\right\rvert}
\providecommand{\dd}{\,\mathrm{d}}
\providecommand{\e}{\mathrm{e}}
\providecommand{\Card}{\operatorname{Card}}
\providecommand{\E}{\operatorname{E}}
\providecommand{\V}{\operatorname{V}}
\providecommand{\degre}{^{\circ}}
\providecommand{\un}[1]{\,\mathrm{#1}}
\providecommand{\sci}[2]{#1\times10^{#2}}
\providecommand{\pgcd}{\operatorname{PGCD}}
\providecommand{\ppcm}{\operatorname{PPCM}}
\title{Prime numbers and Fermat's little theorem}
\subtitle{Lecture and worked examples}
\date{}
\author{}
\begin{document}
\begin{frame}[plain]
\vfill
{\usebeamerfont{title}\inserttitle\par}
\medskip
{\usebeamerfont{subtitle}\insertsubtitle\par}
\medskip
{\usebeamercolor[fg]{progress bar}\rule{\textwidth}{.5pt}}
\vfill
\end{frame}
\begin{frame}{Objectives}
Recognise a prime number; prove that there are infinitely many primes; factorise an integer into primes and use the factorisation (divisors, gcd); use Fermat's little theorem to compute remainders.
\end{frame}
\begin{frame}{Introductory activity}
For an integer \(n\ge2\) and an integer \(a\), let \(r_n(a)\) be the remainder of \(a^{n-1}\) in the division by \(n\).
\end{frame}
\begin{frame}{Introductory activity (2)}
\begin{enumerate}
\setcounter{enumi}{0}
\item Compute \(r_5(2)\), \(r_5(3)\), \(r_7(2)\), \(r_7(3)\) by hand, using congruences.

\item Using the Python function \texttt{pow(a, n - 1, n)}, which returns this remainder, complete the following table.
  
\begin{center}\begin{tabular}{llllllll}
\(n\) & \(5\) & \(6\) & \(7\) & \(8\) & \(9\) & \(10\) & \(11\) \\ 
\(r_n(2)\) &  &  &  &  &  &  &  \\ 
\(r_n(3)\) &  &  &  &  &  &  &  \\ 
\end{tabular}\end{center}

\item For which values of \(n\) do we get \(1\) on both rows? What do they have in common?

\item State a conjecture. Test it for \(n=13\) and \(a=5\).

\item For \(n=7\), compute \(r_7(7)\) and \(r_7(14)\). How should the conjecture be made more precise?

\end{enumerate}
\end{frame}
\begin{frame}{Introductory activity (3)}
\emph{Observing is not proving: the properties conjectured here are established in the course.}
\end{frame}
\begin{frame}{Prime numbers}
A natural number \(p\) is \emph{prime} if it has exactly two positive divisors, \(1\) and \(p\). Thus \(1\) is not prime.
\end{frame}
\begin{frame}{Prime numbers}
Every integer \(n\ge2\) has a prime divisor: its smallest divisor greater than or equal to \(2\). If \(n\) is not prime, this divisor \(p\) satisfies \(p^2\le n\).
\end{frame}
\begin{frame}{Prime numbers}
Let \(p\) be the smallest divisor of \(n\) greater than or equal to \(2\) (it exists, since \(n\) is one). If \(p\) were not prime, it would have a divisor \(d\) with \(2\le d<p\), which would also divide \(n\): a contradiction. If \(n\) is not prime, \(n=pq\) with \(q\ge2\) a divisor of \(n\), so \(q\ge p\) and \(n=pq\ge p^2\).
\end{frame}
\begin{frame}{Prime numbers}
To know whether \(n\ge2\) is prime, it is enough to check that no prime \(p\) with \(p^2\le n\) divides \(n\). Example: \(211\) is divisible neither by \(2\), \(3\), \(5\), \(7\), \(11\), nor by \(13\), and \(17^2=289>211\): \(211\) is prime.
\end{frame}
\begin{frame}{Prime numbers}
There are infinitely many prime numbers.
\end{frame}
\begin{frame}{Prime numbers}
Assume there are only finitely many: \(p_1,\dots,p_k\). Let \(N=p_1p_2\cdots p_k+1\ge2\). By the previous theorem, \(N\) has a prime divisor, which is one of the \(p_i\). This \(p_i\) divides \(N\) and the product \(p_1\cdots p_k\), so it divides their difference \(1\): this is impossible. The assumption is false.
\end{frame}
\begin{frame}{Prime numbers}
Cross out the multiples of each prime starting from its square; the numbers not crossed out are prime.
lstlisting
def crible(n):
    premier = [True] * (n + 1)
    premier[0] = premier[1] = False
    p = 2
    while p * p <= n:
        if premier[p]:
            for m in range(p * p, n + 1, p):
                premier[m] = False
        p = p + 1
    return [k for k in range(n + 1) if premier[k]]
lstlisting
There are \(25\) prime numbers less than \(100\).
\end{frame}
\begin{frame}{Prime numbers}
Let \(p\) be a prime number. If \(p\mid ab\), then \(p\mid a\) or \(p\mid b\). More generally, if \(p\) divides a product, it divides one of the factors.
\end{frame}
\begin{frame}{Prime numbers}
If \(p\nmid a\), then \(\operatorname{gcd}(p,a)=1\), because the only positive divisors of \(p\) are \(1\) and \(p\). Gauss's theorem gives \(p\mid b\). The case of a product of several factors follows by induction on their number.
\end{frame}
\begin{frame}{Prime factorisation}
Every integer \(n\ge2\) can be written as a product of primes, \(n=p_1^{\alpha_1}p_2^{\alpha_2}\cdots p_k^{\alpha_k}\), with \(p_1<p_2<\dots<p_k\) prime and \(\alpha_i\ge1\). This expression is unique.
\end{frame}
\begin{frame}{Prime factorisation}
\emph{Existence}, by strong induction. The integer \(2\) is prime. Assume every integer in \(\{2,\ldots,n\}\) can be factorised. If \(n+1\) is prime, we are done; otherwise \(n+1=ab\) with \(2\le a,b\le n\), and we put the factorisations of \(a\) and \(b\) together.

\emph{Uniqueness.} If \(p_1\cdots p_r=q_1\cdots q_s\) (repeated prime factors), \(p_1\) divides the product of the \(q_j\), hence one of them, which is prime: \(p_1=q_j\). We cancel and start again; the two lists of factors coincide.
\end{frame}
\begin{frame}{Prime factorisation}
Divide by the smallest divisor \(p\ge2\) as long as possible; it is prime. When \(p^2>n\), the remaining quotient, if it exceeds \(1\), is prime.
lstlisting
def decomposition(n):
    facteurs = []
    p = 2
    while p * p <= n:
        while n \% p == 0:
            facteurs.append(p)
            n = n // p
        p = p + 1
    if n > 1:
        facteurs.append(n)
    return facteurs
lstlisting
For example, \texttt{decomposition(360)} returns \texttt{[2, 2, 2, 3, 3, 5]}.
\end{frame}
\begin{frame}{Prime factorisation}
If \(n=p_1^{\alpha_1}\cdots p_k^{\alpha_k}\), the positive divisors of \(n\) are the integers \(p_1^{\beta_1}\cdots p_k^{\beta_k}\) with \(0\le\beta_i\le\alpha_i\). There are \((\alpha_1+1)(\alpha_2+1)\cdots(\alpha_k+1)\) of them.
\end{frame}
\begin{frame}{Prime factorisation}
The gcd of two integers is the product of the common prime factors, each raised to the smaller of the two exponents. Example: \(360=2^3\times3^2\times5\) and \(84=2^2\times3\times7\) give \(\operatorname{gcd}(360,84)=2^2\times3=12\).
\end{frame}
\begin{frame}{Fermat's little theorem}
Let \(p\) be a prime number and \(a\) an integer. Then \(a^p\equiv a\pmod p\). If moreover \(p\) does not divide \(a\), then \(a^{p-1}\equiv1\pmod p\).
\end{frame}
\begin{frame}{Fermat's little theorem}
\emph{Step 1.} For \(1\le k\le p-1\), \(p\) divides \(\binom pk\). Indeed \(k\binom pk=p\binom{p-1}{k-1}\), so \(p\mid k\binom pk\); since \(1\le k<p\), \(p\) does not divide \(k\) and \(\operatorname{gcd}(p,k)=1\); Gauss gives \(p\mid\binom pk\).

\emph{Step 2.} By the binomial formula, for every integer \(a\), \((a+1)^p=a^p+\sum_{k=1}^{p-1}\binom pk a^k+1\equiv a^p+1\pmod p\).
\end{frame}
\begin{frame}{Fermat's little theorem (2)}
\emph{Step 3.} Induction on \(a\ge0\): \(0^p\equiv0\), and if \(a^p\equiv a\), then \((a+1)^p\equiv a^p+1\equiv a+1\pmod p\). For \(a<0\), choose an integer \(b\ge0\) such that \(b\equiv a\pmod p\); then \(a^p\equiv b^p\equiv b\equiv a\).

\emph{Step 4.} If \(p\nmid a\), then \(p\mid a^p-a=a\left(a^{p-1}-1\right)\) and \(\operatorname{gcd}(p,a)=1\); Gauss gives \(p\mid a^{p-1}-1\).
\end{frame}
\begin{frame}{Fermat's little theorem}
Remainder of \(3^{2026}\) in the division by \(11\): \(11\) is prime and does not divide \(3\), so \(3^{10}\equiv1\pmod{11}\). Since \(2{,}026=10\times202+6\), \(3^{2026}\equiv3^6=729\equiv3\pmod{11}\) (because \(729=66\times11+3\)).
\end{frame}
\begin{frame}{Fermat's little theorem}
The converse is false: \(2^{340}\equiv1\pmod{341}\), whereas \(341=11\times31\).
\end{frame}
\end{document}
